\section{Теорема о нормальном базисе}

\begin{definition}[5.17]
Пусть $E/F$~--- конечное расширение Галуа с группой $G$. Базис $E$ как $F$-векторного пространства называется
\emph{нормальным}, если он состоит из сопряжённых одного элемента $E$, т.\,е.\ имеет вид $\{\sigma\alpha\mid\sigma\in G\}$
для некоторого $\alpha\in E$. Такой элемент $\alpha$ будем называть \emph{нормальным}.
\end{definition}

Так как $|G|=[E:F]$, элементы $\sigma\alpha$ ($\sigma\in G$) образуют базис тогда и только тогда, когда они линейно
независимы над $F$ (в частности, попарно различны).

\begin{example}[5.17a]
(a) $E=\Q[\sqrt2]$, $G=\{1,\sigma\}$, $\sigma\sqrt2=-\sqrt2$. Для $\alpha=a+b\sqrt2$ ($a,b\in\Q$) имеем
$\sigma\alpha=a-b\sqrt2$, и координаты $\alpha,\sigma\alpha$ в базисе $1,\sqrt2$ образуют матрицу
$\left(\begin{smallmatrix}a&b\\ a&-b\end{smallmatrix}\right)$ с определителем $-2ab$. Значит, $\alpha$ нормален $\iff$
$ab\ne0$. Например, $\sqrt2$ порождает $E$ как поле, но $\{\sqrt2,-\sqrt2\}$ не базис; а $\{1+\sqrt2,1-\sqrt2\}$~---
нормальный базис. Точно так же в $\Q[i]$ сопряжённые $i,-i$ базиса не образуют, а $1+i,1-i$ образуют.

(b) Пусть $p$ простое, $\zeta$~--- первообразный корень степени $p$ из $1$, $E=\Q[\zeta]$. Группа Галуа состоит из
автоморфизмов $\zeta\mapsto\zeta^k$, $1\le k\le p-1$, так что сопряжённые $\zeta$~--- это $\zeta,\zeta^2,\dots,\zeta^{p-1}$.
Они получаются из базиса $1,\zeta,\dots,\zeta^{p-2}$ умножением на $\zeta$, т.\,е.\ обратимым $\Q$-линейным
отображением; значит, $\{\zeta,\zeta^2,\dots,\zeta^{p-1}\}$~--- нормальный базис.

(c) $\F_4=\F_2[\omega]$, $\omega^2+\omega+1=0$; группа Галуа над $\F_2$ порождена отображением $x\mapsto x^2$, и
$\omega^2=\omega+1$. Элементы $\omega,\omega+1$ линейно независимы над $\F_2$~--- это нормальный базис. Элемент $1$ не
нормален: его орбита состоит из одного элемента.
\end{example}

\begin{theorem}[5.18: теорема о нормальном базисе]
Каждое конечное расширение Галуа обладает нормальным базисом.
\end{theorem}

\subsection*{Групповая алгебра}

\emph{Групповая алгебра} $F[G]$ (обозначается также $FG$) конечной группы $G$ над полем $F$~--- это $F$-векторное
пространство с базисом $\{\sigma\mid\sigma\in G\}$ и умножением, продолжающим умножение в $G$ по линейности:
\[
\Bigl(\sum_{\sigma\in G}a_\sigma\sigma\Bigr)\Bigl(\sum_{\tau\in G}b_\tau\tau\Bigr)
=\sum_{\rho\in G}\Bigl(\sum_{\sigma\tau=\rho}a_\sigma b_\tau\Bigr)\rho,\qquad a_\sigma,b_\tau\in F.
\]
Это ассоциативная $F$-алгебра с единицей $1=1_G$; она коммутативна тогда и только тогда, когда $G$ абелева.

\emph{Универсальное свойство.} Если $G$ действует $F$-линейно на $F$-векторном пространстве $V$, то это действие
единственным образом продолжается до структуры $F[G]$-модуля: $\bigl(\sum a_\sigma\sigma\bigr)v=\sum a_\sigma\,\sigma(v)$.
Обратно, $F[G]$-модуль~--- это $F$-пространство с линейным действием $G$. Гомоморфизмы $F[G]$-модулей~--- это
$F$-линейные отображения, перестановочные с действием $G$.

Пусть $E/F$~--- конечное расширение Галуа с группой $G$. Элементы $G$ действуют на $E$ $F$-линейно, поэтому $E$~---
$F[G]$-модуль: $\bigl(\sum a_\sigma\sigma\bigr)x=\sum a_\sigma\,\sigma(x)$.

\begin{proposition}[5.18a]
Для $\alpha\in E$ отображение
\[
\varphi_\alpha\colon F[G]\to E,\qquad \sum_\sigma a_\sigma\sigma\mapsto\sum_\sigma a_\sigma\,\sigma\alpha,
\]
является гомоморфизмом $F[G]$-модулей, и $\{\sigma\alpha\}_{\sigma\in G}$~--- нормальный базис тогда и только тогда, когда
$\varphi_\alpha$~--- изоморфизм. Поэтому теорема \rf{5.18} равносильна тому, что \emph{$E$~--- свободный $F[G]$-модуль
ранга~$1$}: $E\cong F[G]$ как $F[G]$-модули.
\end{proposition}

\begin{proof}
$\varphi_\alpha$ $F$-линейно и $\varphi_\alpha(\tau x)=\tau\varphi_\alpha(x)$ для $\tau\in G$:
$\varphi_\alpha\bigl(\sum a_\sigma\tau\sigma\bigr)=\sum a_\sigma\tau\sigma\alpha=\tau\bigl(\sum a_\sigma\sigma\alpha\bigr)$.
Отображение $\varphi_\alpha$ переводит базис $\{\sigma\}$ пространства $F[G]$ в семейство $\{\sigma\alpha\}$, и оно биективно
тогда и только тогда, когда это семейство~--- базис $E$. Наконец, если $\psi\colon F[G]\to E$~--- какой-нибудь изоморфизм
$F[G]$-модулей и $\alpha=\psi(1)$, то $\psi(\sum a_\sigma\sigma)=\sum a_\sigma\sigma\psi(1)$, т.\,е.\ $\psi=\varphi_\alpha$.
\end{proof}

Мы приведём три доказательства теоремы \rf{5.18}: для бесконечного $F$; для циклической группы $G$ (это покрывает
конечные поля, так как расширения Галуа конечного поля циклические, \rf{4.20}, \rf{4.23}); и равномерное, использующее теорему
Крулля--Шмидта.

\subsection*{Доказательство для бесконечного поля}

\begin{lemma}[5.19]
Пусть $f\in F[X_1,\dots,X_m]$ и $S$~--- бесконечное подмножество $F$. Если $f(a_1,\dots,a_m)=0$ для всех
$a_1,\dots,a_m\in S$, то $f$~--- нулевой многочлен.
\end{lemma}

\begin{proof}
Индукция по $m$. При $m=1$ ненулевой многочлен от одной переменной имеет не больше $\deg f$ корней (\rf{1.7}), а у $f$
их бесконечно много. Пусть $m>1$; запишем $f$ как многочлен от $X_m$ с коэффициентами из $F[X_1,\dots,X_{m-1}]$:
\[
f=\sum_ic_i(X_1,\dots,X_{m-1})\,X_m^i.
\]
Для любых $a_1,\dots,a_{m-1}\in S$ многочлен $f(a_1,\dots,a_{m-1},X_m)=\sum_ic_i(a_1,\dots,a_{m-1})X_m^i\in F[X_m]$
обращается в нуль во всех точках бесконечного множества $S$, поэтому (случай $m=1$) все его коэффициенты
$c_i(a_1,\dots,a_{m-1})$ равны нулю. Это верно для всех $a_1,\dots,a_{m-1}\in S$, и по предположению индукции все
$c_i$~--- нулевые многочлены. Значит, $f=0$.
\end{proof}

Лемма неверна для конечного $S$: многочлен $X^q-X\ne0$ обращается в нуль во всех точках $\F_q$. Именно поэтому
первое доказательство работает только для бесконечных полей.

\begin{proof}[Доказательство теоремы \rf{5.18} для бесконечного $F$]
Пусть $m=|G|=[E:F]$; занумеруем элементы группы, начиная с единицы: $G=\{\sigma_1,\dots,\sigma_m\}$, ${\sigma_1=1}$.

\emph{Шаг 1: если $f\in E[X_1,\dots,X_m]$ и $f(\sigma_1\alpha,\dots,\sigma_m\alpha)=0$ для всех $\alpha\in E$, то
$f=0$.} Выберем базис $\alpha_1,\dots,\alpha_m$ поля $E$ над $F$ и положим $B=(\sigma_j\alpha_i)_{j,i}$; по \rf{5.16} эта
матрица обратима. Пусть $g\in E[Y_1,\dots,Y_m]$ получен из $f$ линейной заменой $X_j=\sum_iB_{ji}Y_i$:
\[
g(Y_1,\dots,Y_m)=f\Bigl(\sum_{i=1}^mY_i\,\sigma_1\alpha_i,\ \dots,\ \sum_{i=1}^mY_i\,\sigma_m\alpha_i\Bigr).
\]
Для $a_1,\dots,a_m\in F$ элемент $\alpha=\sum_ia_i\alpha_i$ пробегает всё $E$, и, так как $\sigma_j$ $F$-линейны,
\[
g(a_1,\dots,a_m)=f\Bigl(\sigma_1\Bigl(\sum_ia_i\alpha_i\Bigr),\dots,\sigma_m\Bigl(\sum_ia_i\alpha_i\Bigr)\Bigr)
=f(\sigma_1\alpha,\dots,\sigma_m\alpha)=0.
\]
Применим лемму \rf{5.19} к полю $E$ и его бесконечному подмножеству $S=F$: получаем $g=0$. Замены $X\mapsto BY$ и
$Y\mapsto B^{-1}X$ задают взаимно обратные гомоморфизмы $E$-алгебр $E[X_1,\dots,X_m]\rightleftarrows E[Y_1,\dots,Y_m]$
(на образующих их композиции тождественны), поэтому $f(X)=g(B^{-1}X)=0$. Здесь и только здесь используется
бесконечность $F$.

\emph{Шаг 2: матрица таблицы умножения.} Положим $X(\sigma_k)=X_k$ и рассмотрим матрицу $A=\bigl(X(\sigma_i\sigma_j)\bigr)_{i,j}$:
на месте $(i,j)$ стоит $X_k$, где $\sigma_i\sigma_j=\sigma_k$. Это таблица умножения (таблица Кэли) группы $G$, в которой
каждый элемент $\sigma_k$ заменён переменной $X_k$. Пусть $h(X_1,\dots,X_m)=\det A$; это однородный многочлен
степени $m$ с целыми коэффициентами (точнее, с их образами в простом подполе $F$). При $X_1=1$, $X_2=\dots=X_m=0$ на месте
$(i,j)$ стоит $1$ тогда и только тогда, когда $\sigma_i\sigma_j=1$, т.\,е.\ $\sigma_j=\sigma_i^{-1}$; в каждой строке и
каждом столбце ровно одна единица. Это матрица перестановки, поэтому $h(1,0,\dots,0)=\pm1\ne0$, и $h$~--- ненулевой
многочлен.

\emph{Шаг 3.} По шагу 1 найдётся $\alpha\in E$ с
\[
h(\sigma_1\alpha,\dots,\sigma_m\alpha)=\det(\sigma_i\sigma_j\alpha)_{i,j}\ne0.
\]
Докажем, что $\sigma_1\alpha,\dots,\sigma_m\alpha$~--- нормальный базис; так как их $m=[E:F]$, достаточно линейной
независимости. Пусть $\sum_ja_j\,\sigma_j\alpha=0$, $a_j\in F$. Применяя $\sigma_1,\dots,\sigma_m$ (они тождественны на
$F$), получаем систему из $m$ уравнений
\[
\sum_{j=1}^m(\sigma_i\sigma_j\alpha)\,a_j=0,\qquad i=1,\dots,m,
\]
с обратимой матрицей $(\sigma_i\sigma_j\alpha)$. Значит, все $a_j=0$.
\end{proof}

\begin{center}
\begin{tabular}{c|ccc}
$\cdot$ & $1$ & $\sigma$ & $\sigma^2$\\ \hline
$1$ & $1$ & $\sigma$ & $\sigma^2$\\
$\sigma$ & $\sigma$ & $\sigma^2$ & $1$\\
$\sigma^2$ & $\sigma^2$ & $1$ & $\sigma$
\end{tabular}
\qquad$\leadsto$\qquad
$A=\begin{pmatrix}\color{cnew}X_1&X_2&X_3\\ X_2&X_3&\color{cnew}X_1\\ X_3&\color{cnew}X_1&X_2\end{pmatrix}$,
\qquad
$h(1,0,0)=\det\begin{pmatrix}\color{cnew}1&0&0\\0&0&\color{cnew}1\\0&\color{cnew}1&0\end{pmatrix}=-1.$
\end{center}
\figcap{Шаг 2 для $G=\{1,\sigma,\sigma^2\}$ ($\sigma_1=1$, $\sigma_2=\sigma$, $\sigma_3=\sigma^2$): таблица Кэли
превращается в матрицу $A$; единицы матрицы перестановки стоят там, где $\sigma_i\sigma_j=1$.}

\begin{example}[5.19a]
Для группы порядка $3$ из рисунка $h=\det A=-(X_1^3+X_2^3+X_3^3-3X_1X_2X_3)
=-(X_1+X_2+X_3)(X_1^2+X_2^2+X_3^2-X_1X_2-X_1X_3-X_2X_3)$. Для группы порядка $2$ имеем
$A=\left(\begin{smallmatrix}X_1&X_2\\X_2&X_1\end{smallmatrix}\right)$ и $h=X_1^2-X_2^2$. В примере \rf{5.17a}(a), где
$\alpha=a+b\sqrt2$, получаем $h(\alpha,\sigma\alpha)=\alpha^2-(\sigma\alpha)^2=4ab\sqrt2$: условие шага~3 снова
даёт $ab\ne0$. (С точностью до знака $h$~--- это \emph{групповой определитель} группы $G$; столбцы матрицы
$\bigl(X(\sigma_i\sigma_j^{-1})\bigr)$ получаются из столбцов $A$ перестановкой.)
\end{example}

\subsection*{Доказательство для циклической группы}

Пусть $V$~--- конечномерное пространство над $F$, $T\colon V\to V$~--- линейный оператор. Для $Q\in F[X]$ и $v\in V$
положим $Q\cdot v=Q(T)v$; так $V$ становится $F[X]$-модулем. \emph{Минимальный многочлен} $P$ оператора $T$~---
унитарный многочлен наименьшей степени с $P(T)=0$. Для $v\in V$ обозначим через $P_v$ унитарный многочлен наименьшей
степени с $P_v(T)v=0$ (он существует, так как $P(T)v=0$). Деление с остатком показывает: \emph{если $Q(T)v=0$, то
$P_v\mid Q$} (остаток $R=Q-SP_v$ тоже обнуляет $v$ и имеет меньшую степень, значит, $R=0$). Аналогично $P$ делит
любой $Q$ с $Q(T)=0$. В частности, $P_v\mid P$.

\begin{lemma}[5.19b: циклический вектор]
Пусть $\dim_FV=n$ и $P$~--- минимальный многочлен оператора $T$.
\begin{enumerate}[label=\textup{(\alph*)}]
\item Существует $v\in V$ с $P_v=P$.
\item Если $\deg P=n$, то для такого $v$ векторы $v,Tv,\dots,T^{n-1}v$ образуют базис $V$, и отображение $Q\mapsto Q(T)v$
задаёт изоморфизм $F[X]$-модулей $F[X]/(P)\to V$.
\end{enumerate}
\end{lemma}

\begin{proof}
(i) \emph{Если $R$ унитарен и $R\mid P_u$, то $P_{R(T)u}=P_u/R$.} Действительно, $Q(T)R(T)u=0\iff P_u\mid QR\iff
(P_u/R)\mid Q$.

(ii) \emph{Если $P_u$ и $P_w$ взаимно просты, то $P_{u+w}=P_uP_w$.} Ясно, что $(P_uP_w)(T)$ обнуляет $u+w$. Пусть
$Q(T)(u+w)=0$, т.\,е.\ $Q(T)u=-Q(T)w$. Применив $P_w(T)$, получим $P_w(T)Q(T)u=-Q(T)P_w(T)w=0$, откуда $P_u\mid P_wQ$ и,
по взаимной простоте, $P_u\mid Q$. Так же $P_w\mid Q$, и $P_uP_w\mid Q$.

(iii) \emph{Для любых $u,w\in V$ найдётся $z$ с $P_z=\operatorname{\textup{НОК}}(P_u,P_w)$.} Разложим на неприводимые:
$P_u=\prod_qq^{a_q}$, $P_w=\prod_qq^{b_q}$. Положим $A=\prod_{a_q\ge b_q}q^{a_q}$, $B=\prod_{a_q<b_q}q^{b_q}$. Тогда
$A\mid P_u$, $B\mid P_w$, $A$ и $B$ взаимно просты и $AB=\operatorname{\textup{НОК}}(P_u,P_w)$. По (i) у векторов
$u'=(P_u/A)(T)u$ и $w'=(P_w/B)(T)w$ минимальные многочлены $A$ и $B$, а по (ii) $z=u'+w'$ годится.

(a) Пусть $e_1,\dots,e_n$~--- базис $V$ и $L=\operatorname{\textup{НОК}}(P_{e_1},\dots,P_{e_n})$. По (iii) и индукции найдётся $v$
с $P_v=L$. Так как $L(T)$ обнуляет базис, $L(T)=0$ и $P\mid L$; а так как каждое $P_{e_i}$ делит $P$, то $L\mid P$.
Итак, $P_v=L=P$.

(b) Если $\sum_{i<n}c_iT^iv=0$ и не все $c_i$ нулевые, то $Q=\sum_{i<n}c_iX^i$~--- ненулевой многочлен степени $<n$ с
$Q(T)v=0$, и $P_v\mid Q$~--- невозможно, так как $\deg P_v=n$. Значит, $n$ векторов $v,\dots,T^{n-1}v$ линейно независимы и
образуют базис. Отображение $F[X]\to V$, $Q\mapsto Q(T)v$, $F[X]$-линейно, сюръективно (образ содержит базис), а его ядро
$\{Q\mid Q(T)v=0\}=(P_v)=(P)$.
\end{proof}

\begin{remark}[5.19c]
На языке модулей: $V$~--- конечно порождённый $F[X]$-модуль кручения, и по структурной теореме
$V\cong F[X]/(P_1)\oplus\dots\oplus F[X]/(P_r)$ с унитарными непостоянными $P_1\mid P_2\mid\dots\mid P_r$ (инвариантные
множители). Минимальный многочлен $T$ равен $P_r$ (наименьшее общее кратное многочленов $P_i$), а характеристический
многочлен равен $P_1\cdots P_r$ (характеристический многочлен умножения на $X$ в $F[X]/(Q)$~--- это $Q$: в базисе
$1,X,\dots,X^{\deg Q-1}$ это сопровождающая матрица $Q$). Поэтому \emph{минимальный многочлен совпадает с
характеристическим $\iff r=1\iff V\cong F[X]/(P)$~--- циклический модуль}. По теореме Гамильтона--Кэли $P$ делит
характеристический многочлен степени $n=\dim V$, так что условие $\deg P=n$ в \rf{5.19b}(b) означает в точности совпадение
минимального и характеристического многочленов. Лемма \rf{5.19b} доказывает нужную часть структурной теоремы напрямую.

Для бесконечного $F$ пункт (a) леммы \rf{5.19b} получается короче. Если $P_v\ne P$ для всех $v$, то каждый $v$ лежит в
$\ker Q(T)$ для некоторого унитарного собственного делителя $Q$ многочлена $P$ (а именно $Q=P_v$); таких $Q$ конечное
число, и каждое $\ker Q(T)$~--- собственное подпространство ($Q(T)\ne0$ по минимальности $P$). Но над бесконечным полем
пространство не является объединением конечного числа собственных подпространств $V_1,\dots,V_k$: выберем ненулевые
линейные функции $\ell_i$ с $V_i\subset\ker\ell_i$; произведение $\ell_1\cdots\ell_k$~--- ненулевой многочлен от
координат (кольцо многочленов~--- область целостности), обращающийся в нуль на всём $V\cong F^n$, что противоречит лемме
\rf{5.19}. Над конечным полем это рассуждение не проходит, и там нужны пункты (i)--(iii).
\end{remark}

\begin{proof}[Доказательство теоремы \rf{5.18} для циклической $G$]
Пусть $G=\langle\sigma\rangle$ имеет порядок $n=[E:F]$, и $P$~--- минимальный многочлен $\sigma$ как $F$-линейного
оператора на $E$. Так как $\sigma^n=1$, то $P\mid X^n-1$. С другой стороны, $1,\sigma,\dots,\sigma^{n-1}$~--- различные
автоморфизмы $E$, и по \rf{5.15} они линейно независимы над $E$, а тем более над $F$: если $Q=\sum_{i<n}c_iX^i\in F[X]$ и
$Q(\sigma)=\sum c_i\sigma^i=0$, то все $c_i=0$. Поэтому $\deg P\ge n$ и $P=X^n-1$. По лемме \rf{5.19b} найдётся
$\alpha\in E$, для которого $\alpha,\sigma\alpha,\dots,\sigma^{n-1}\alpha$~--- базис $E$; это нормальный базис, так как
$G=\{1,\sigma,\dots,\sigma^{n-1}\}$.
\end{proof}

Итак, $E\cong F[X]/(X^n-1)$ как $F[X]$-модули, где $X$ действует как $\sigma$, а $\alpha$ соответствует классу $1$; базис
$1,X,\dots,X^{n-1}$ переходит в $\alpha,\sigma\alpha,\dots,\sigma^{n-1}\alpha$. Так как $F[G]\cong F[X]/(X^n-1)$
($\sigma\leftrightarrow X$), это и есть изоморфизм $F[G]\cong E$ из \rf{5.18a}. Обратите внимание, что сепарабельность
$X^n-1$ не нужна: при $\operatorname{char}F=p\mid n$ многочлен $X^n-1$ имеет кратные множители и $\sigma$ не
диагонализуем, но лемма \rf{5.19b} этого не требует.

\begin{remark}[5.19d]
Для конечного $F=\F_q$ теорему Дедекинда можно заменить подсчётом. Здесь $E=\F_{q^n}$ и $\sigma(x)=x^q$
(\rf{4.20}, \rf{4.23}). Если $Q=\sum_{i=0}^dc_iX^i\ne0$, $c_d\ne0$, $d<n$, то $Q(\sigma)(x)=\sum_ic_ix^{q^i}$~--- ненулевой
многочлен от $x$ степени $q^d<q^n=|E|$, и он не может обращаться в нуль на всём $E$. Значит, $\deg P\ge n$.
\end{remark}

\begin{example}[5.19e]
Пусть $E=\F_8=\F_2[x]/(x^3+x+1)$, $\sigma(y)=y^2$. Элемент $x$ не нормален: его сопряжённые $x$, $x^2$, $x^4=x^2+x$ в
сумме дают $0$. Элемент $x+1$ нормален: его сопряжённые $x+1$, $x^2+1$, $x^4+1=x^2+x+1$ имеют в базисе $1,x,x^2$
координаты $(1,1,0)$, $(1,0,1)$, $(1,1,1)$, и определитель этой матрицы равен $1$.

Все нормальные элементы описываются так. Пусть $R=F[X]/(X^n-1)$ и $\alpha$ нормален, т.\,е.\ $r\mapsto r\alpha$~---
изоморфизм $R\to E$. Элемент $u\alpha$ ($u\in R$) нормален $\iff$ отображение $r\mapsto ru\alpha$ биективно $\iff$
умножение на $u$ биективно на $R$ $\iff$ $u\in R^\times$. Здесь $X^3-1=(X+1)(X^2+X+1)$ над $\F_2$, и по китайской
теореме об остатках $R\cong\F_2\times\F_4$, так что $|R^\times|=1\cdot3=3$. Значит, в $\F_8$ ровно три нормальных
элемента: $x+1$, $x^2+1$, $x^2+x+1$ (одна орбита Фробениуса); это легко проверить и прямым перебором семи ненулевых
элементов.
\end{example}

\subsection*{Равномерное доказательство}

Модуль над кольцом называется \emph{неразложимым}, если он ненулевой и не является прямой суммой двух ненулевых
подмодулей. Модуль имеет \emph{конечную длину}, если длины цепочек его подмодулей $M_0\subsetneq M_1\subsetneq\dots$
ограничены. Например, $F[G]$-модуль, конечномерный над $F$, имеет конечную длину: его подмодули~--- подпространства.

\begin{theorem}[5.19f: Крулль--Шмидт]
Каждый ненулевой модуль $M$ конечной длины над кольцом раскладывается в прямую сумму неразложимых модулей, и
неразложимые слагаемые определены однозначно с точностью до порядка и изоморфизма: если
$M\cong M_1\oplus\dots\oplus M_r\cong N_1\oplus\dots\oplus N_s$ с неразложимыми $M_i,N_j$, то $r=s$ и после перенумерации
$M_i\cong N_i$.
\end{theorem}

Существование очевидно (индукция по длине: неразложимый модуль раскладывать не нужно, а иначе $M=M'\oplus M''$ со
слагаемыми меньшей длины); доказательство единственности~--- в отступлении \rf{5.19j} ниже. Сгруппировав
изоморфные слагаемые, получаем $M\cong\bigoplus_im_iM_i$, где $M_i$ попарно неизоморфны, а $m_iM_i$~--- сумма $m_i$
копий $M_i$; классы изоморфизма $M_i$ и кратности $m_i$ определены однозначно.

\begin{corollary}[5.19g]
Если $M$ и $M'$~--- модули конечной длины и $M^k\cong M'^k$ для некоторого $k\ge1$, то $M\cong M'$.
\end{corollary}

\begin{proof}
Если один из модулей нулевой, то и другой нулевой. Иначе запишем $M\cong\bigoplus_im_iM_i$,
$M'\cong\bigoplus_im_i'M_i$, где $M_i$~--- попарно неизоморфные неразложимые модули, входящие хотя бы в одно из разложений
(кратности $m_i,m'_i\ge0$). Тогда $M^k\cong\bigoplus_ikm_iM_i$ и $M'^k\cong\bigoplus_ikm_i'M_i$, и по \rf{5.19f}
$km_i=km_i'$, т.\,е.\ $m_i=m_i'$ для всех $i$. Значит, $M\cong M'$.
\end{proof}

Рассмотрим тензорное произведение $E\otimes_FE$. Пусть $E$ действует на первом множителе, а $G$~--- на втором:
\[
a(x\otimes y)=ax\otimes y\quad(a\in E),\qquad \sigma(x\otimes y)=x\otimes\sigma y\quad(\sigma\in G).
\]
Оба действия корректны (отображения $(x,y)\mapsto ax\otimes y$ и $(x,y)\mapsto x\otimes\sigma y$ $F$-билинейны) и
перестановочны; так $E\otimes_FE$ становится $E$-векторным пространством и $F[G]$-модулем.

\begin{lemma}[5.19h]
Пусть $E/F$~--- расширение степени $n$ с базисом $\alpha_1,\dots,\alpha_n$. Тогда $1\otimes\alpha_1,\dots,1\otimes\alpha_n$~---
базис $E\otimes_FE$ над $E$; в частности, $\dim_E(E\otimes_FE)=n$ и $\dim_E\Hom_E(E\otimes_FE,E)=n$.
\end{lemma}

\begin{proof}
Воспользуемся тем, что $\alpha_j\otimes\alpha_i$ ($1\le i,j\le n$)~--- базис $E\otimes_FE$ над $F$. Элементы
$1\otimes\alpha_i$ порождают $E\otimes_FE$ над $E$: $x\otimes y=\sum_iy_i\,x(1\otimes\alpha_i)$, если
$y=\sum_iy_i\alpha_i$, $y_i\in F$. Пусть $\sum_ic_i(1\otimes\alpha_i)=0$, $c_i\in E$; запишем $c_i=\sum_jc_{ij}\alpha_j$,
$c_{ij}\in F$. Тогда
\[
0=\sum_ic_i\otimes\alpha_i=\sum_{i,j}c_{ij}\,\alpha_j\otimes\alpha_i,
\]
откуда все $c_{ij}=0$ и все $c_i=0$. Пространство $E$-линейных функционалов на $n$-мерном пространстве тоже $n$-мерно.
\end{proof}

\begin{proof}[Равномерное доказательство теоремы \rf{5.18}]
Пусть $E/F$~--- расширение Галуа с группой $G$, $n=[E:F]=|G|$. Докажем, что
\[
\underbrace{F[G]\oplus\dots\oplus F[G]}_n\;\cong\;E\otimes_FE\;\cong\;\underbrace{E\oplus\dots\oplus E}_n
\qquad\text{как $F[G]$-модули;}
\]
тогда $E\cong F[G]$ по \rf{5.19g}, и по \rf{5.18a} у $E$ есть нормальный базис.

(1) Для $\sigma\in G$ пусть $\lambda_\sigma\colon E\otimes_FE\to E$, $x\otimes y\mapsto x\cdot\sigma y$ (корректно, так как
$(x,y)\mapsto x\sigma y$ $F$-билинейно). Ясно, что $\lambda_\sigma$ $E$-линейно, и $\lambda_\sigma(\tau z)=\lambda_{\sigma\tau}(z)$
для $\tau\in G$: на образующих $\lambda_\sigma(x\otimes\tau y)=x\cdot\sigma\tau y$. \emph{Функционалы $\lambda_\sigma$
образуют базис $\Hom_E(E\otimes_FE,E)$.} По \rf{5.19h} достаточно линейной независимости: если $\sum_\sigma c_\sigma\lambda_\sigma=0$,
$c_\sigma\in E$, то $0=\sum_\sigma c_\sigma\lambda_\sigma(1\otimes y)=\sum_\sigma c_\sigma\,\sigma y$ для всех $y\in E$, и по
\rf{5.15} все $c_\sigma=0$.

(2) Пусть $E[G]$~--- групповая алгебра $G$ над $E$, на которой $G$ действует умножением слева. Рассмотрим
\[
\varphi\colon E\otimes_FE\to E[G],\qquad z\mapsto\sum_{\sigma\in G}\lambda_\sigma(z)\,\sigma^{-1}.
\]
Оно $E$-линейно. Если $\varphi(z)=0$, то $\lambda_\sigma(z)=0$ для всех $\sigma$, а так как $\lambda_\sigma$ образуют базис
двойственного пространства, $z=0$. Значит, $\varphi$ инъективно, и, поскольку оба пространства $n$-мерны над $E$,
$\varphi$~--- изоморфизм. Он перестановочен с $G$: подставляя $\rho=\sigma\tau$, т.\,е.\ $\sigma^{-1}=\tau\rho^{-1}$, получаем
\[
\varphi(\tau z)=\sum_\sigma\lambda_{\sigma\tau}(z)\,\sigma^{-1}=\sum_\rho\lambda_\rho(z)\,\tau\rho^{-1}=\tau\varphi(z).
\]
Итак, $\varphi$~--- изоморфизм $F[G]$-модулей. Далее, если $e_1,\dots,e_n$~--- базис $E$ над $F$, то
$E[G]=e_1F[G]\oplus\dots\oplus e_nF[G]$ (запишем коэффициенты $c_\sigma\in E$ элемента $\sum c_\sigma\sigma$ в базисе
$e_i$), каждое слагаемое устойчиво относительно умножения на $G$ слева и изоморфно $F[G]$. Значит,
$E\otimes_FE\cong E[G]\cong F[G]^n$.

(3) С другой стороны, каждый элемент $E\otimes_FE$ однозначно записывается в виде $\sum_ie_i\otimes y_i$, $y_i\in E$,
т.\,е.\ $E\otimes_FE=(e_1\otimes E)\oplus\dots\oplus(e_n\otimes E)$. Каждое слагаемое устойчиво относительно $G$, и
$y\mapsto e_i\otimes y$~--- изоморфизм $F[G]$-модулей $E\to e_i\otimes E$. Значит, $E\otimes_FE\cong E^n$.
\end{proof}

\[
\begin{tikzcd}[column sep=large]
F[G]^n\arrow[r,"\cong"] & E[G] & E\otimes_FE\arrow[l,cnew,"\varphi"',"\cong"]\arrow[r,"\cong"] & E^n
\end{tikzcd}
\]
\figcap{Равномерное доказательство: изоморфизм $\varphi$ строится по функционалам $\lambda_\sigma$; из
$F[G]^n\cong E^n$ теорема Крулля--Шмидта даёт $F[G]\cong E$.}

Для конечных полей теорема о нормальном базисе была высказана Эйзенштейном (1850) и доказана Гензелем (1888).
Э.~Нётер (1932) доказала её для некоторых бесконечных полей, а Дойринг в том же году дал равномерное доказательство;
приведённое выше доказательство~--- его упрощение.

\begin{remark}[5.19i]
Отображение $\Lambda=(\lambda_\sigma)_{\sigma\in G}\colon E\otimes_FE\to\prod_{\sigma\in G}E$~--- изоморфизм не только
$E$-пространств (это доказано в (2)), но и колец: $\lambda_\sigma\bigl((x\otimes y)(x'\otimes y')\bigr)=xx'\sigma(yy')=
\lambda_\sigma(x\otimes y)\lambda_\sigma(x'\otimes y')$. Итак, для расширения Галуа $E\otimes_FE\cong E\times\dots\times E$
($n$ копий): после расширения скаляров до $E$ поле $E$ распадается в произведение $n$ копий $E$. Для несепарабельного
расширения так не бывает: если $F=\F_p(t)$, $E=F[\theta]$, $\theta^p=t$, то элемент $u=\theta\otimes1-1\otimes\theta$
ненулевой (это разность двух различных элементов базиса $\theta^j\otimes\theta^i$), но
$u^p=\theta^p\otimes1-1\otimes\theta^p=t(1\otimes1)-t(1\otimes1)=0$; в произведении полей ненулевых нильпотентов нет.
\end{remark}

\begin{aside}[5.19j: о теореме Крулля--Шмидта]
Ключ к единственности~--- \emph{лемма Фиттинга}: если $M$ имеет конечную длину и $f\in\operatorname{End}(M)$, то
$M=\ker f^k\oplus\operatorname{im}f^k$ при достаточно большом $k$. Действительно, цепочки $\ker f\subset\ker f^2\subset\dots$
и $\operatorname{im}f\supset\operatorname{im}f^2\supset\dots$ стабилизируются, скажем, с номера $k$. Если
$x=f^k(y)\in\ker f^k$, то $y\in\ker f^{2k}=\ker f^k$ и $x=0$; а для любого $x$ имеем $f^kx=f^{2k}y$ для некоторого $y$,
и $x=(x-f^ky)+f^ky$ с $x-f^ky\in\ker f^k$. Отсюда для \emph{неразложимого} $M$ каждый эндоморфизм либо нильпотентен, либо
автоморфизм, и сумма двух неавтоморфизмов~--- неавтоморфизм (а по индукции~--- и любая конечная сумма): если $f,g$
нильпотентны, а $u=f+g$ обратим, то
$u^{-1}g=1-u^{-1}f$, где $u^{-1}f$ не обратим, значит, нильпотентен, и $1-u^{-1}f$ обратим (обратный
$\sum_i(u^{-1}f)^i$), так что $g$ обратим~--- противоречие. Теперь пусть $M=M_1\oplus\dots\oplus M_r=N_1\oplus\dots\oplus N_s$,
$\iota_1,p_1$~--- вложение и проекция для $M_1$, а $\iota'_j,p'_j$~--- для $N_j$. Тогда $1_{M_1}=\sum_jp_1\iota'_jp'_j\iota_1$,
и одно из слагаемых~--- автоморфизм $M_1$. Если $g\circ f$~--- автоморфизм $M_1$, где $f=p'_j\iota_1\colon M_1\to N_j$,
$g=p_1\iota'_j$, то $e=f(gf)^{-1}g$~--- ненулевой идемпотент в $\operatorname{End}(N_j)$, и по неразложимости $N_j$ имеем
$e=1$, откуда $f$~--- изоморфизм $M_1\cong N_j$.

Покажем, что $M=M_1\oplus\bigoplus_{i\ne j}N_i$. Пусть $\pi=\iota'_jp'_j$~--- проекция $M$ на $N_j$ вдоль
$\bigoplus_{i\ne j}N_i$; на $M_1$ она совпадает с изоморфизмом $f\colon M_1\to N_j$. Поэтому $M_1\cap\ker\pi=0$, а для
любого $x\in M$ элемент $y=\iota_1f^{-1}p'_j(x)\in M_1$ удовлетворяет $\pi(y)=\pi(x)$, т.\,е.\
$x-y\in\ker\pi=\bigoplus_{i\ne j}N_i$. Итак, $M=M_1\oplus\bigoplus_{i\ne j}N_i$, и
$M_2\oplus\dots\oplus M_r\cong M/M_1\cong\bigoplus_{i\ne j}N_i$. Индукция по $r$ даёт $r-1=s-1$ и изоморфность остальных
слагаемых после перенумерации.
\end{aside}
